\documentclass[a4paper,11pt]{article}
\usepackage[margin=2.2cm]{geometry}
\usepackage[T1]{fontenc}
\usepackage[utf8]{inputenc}
\usepackage{array}
\usepackage{booktabs}
\usepackage{tabularx}
\usepackage{hyperref}

\setlength{\parindent}{0pt}
\setlength{\parskip}{0.6em}
\renewcommand{\arraystretch}{1.15}
\hypersetup{colorlinks=true,linkcolor=blue,urlcolor=blue}

\newcolumntype{Y}{>{\raggedright\arraybackslash}X}
\newcommand{\term}[1]{\texttt{\detokenize{#1}}}
\newcommand{\cmd}[1]{\texttt{\detokenize{#1}}}
\newcommand{\mainsketch}{4\_4\_rc\_kontrol.ino}

\title{\mainsketch{} Kisa Kullanim Kilavuzu}
\author{}
\date{\today}

\begin{document}
\maketitle
\tableofcontents
\newpage

\section{Amac}
Bu dokuman sadece \mainsketch{} icin guncel davranisi kisa ve pratik bicimde anlatir.
Odak noktasi:
\begin{itemize}
  \item hangi komutlarin gecerli oldugu
  \item seri hatta surekli gelen satirlarin anlami
  \item RC ve serial modun nasil secildigi
\end{itemize}

\section{Kullanilan Dosya}
Bu belge yalnizca \mainsketch{} icin hazirlanmistir.

\section{PWM Pin Eslesmesi}
Guncel motor PWM pin eslesmesi:
\begin{itemize}
  \item FL: \term{LPWM=2}, \term{RPWM=3}
  \item FR: \term{LPWM=10}, \term{RPWM=9}
  \item RL: \term{LPWM=6}, \term{RPWM=5}
  \item RR: \term{LPWM=11}, \term{RPWM=12}
\end{itemize}

\section{Ana Calisma Mantigi}
\begin{itemize}
  \item CH1 direksiyon kanali olarak kullanilir.
  \item CH2 enable / mode kanali olarak kullanilir.
  \item CH3 gaz kanali olarak kullanilir.
  \item Varsayilan surus kanali \term{CH3}, enable kanali \term{CH2}'dir.
  \item CH2 aktif ve arm tamamlanmis ise robot RC modunda surulur.
  \item CH2 pasifken serial pair PWM veya tek teker RPM komutlari kullanilabilir.
  \item CH2 pasif ve aktif serial komut yoksa robot \term{BEKLEME} modundadir.
  \item Iki durumda da encoder RPM okunur.
  \item Ileri ve geri suruste hiz tavani maksimum PWM'in \%60'i ile sinirlidir.
  \item Yumusak durus vardir; ileri/geri ve throttle ile donuste cikislar PWM adimlariyla yavaslar.
  \item Teker hiz dengeleme vardir.
  \item Throttle ile giderken sag ve sol hiz farki kullanilir; agresif miks ile ic taraf daha sert duser ve yuksek steering'de hafif ters yone de gecilebilir.
  \item Ileri giderken bir anda geri komut gelirse, disli korumasi icin kontrollu yon degisimi vardir.
  \item Seri telemetride hem teker RPM'leri hem de teker/robot lineer hizlari verilir.
\end{itemize}

\section{Seri Port Ayarlari}
Bu dosya icin seri ayarlar:
\begin{itemize}
  \item Baud: \term{115200}
  \item Data bits: \term{8}
  \item Parity: \term{None}
  \item Stop bit: \term{1}
  \item Line ending: \term{Newline} veya \term{Both NL \& CR}
\end{itemize}

\section{Hiz Hesabi}
Bu kilavuzdaki lineer hiz alanlari asagidaki varsayimla hesaplanir:
\begin{itemize}
  \item Teker yaricapi: \term{19.0 cm}
  \item Teker cevresi: yaklasik \term{119.38 cm}
  \item Teker hizi: \term{RPM x cevre / 60}
\end{itemize}

Bu nedenle RPM degismez; yalnizca \term{cm/s} cinsinden hiz alanlari bu yaricap bilgisine gore hesaplanir.

\section{Acilis Kalibrasyonu}
Kart acildiginda RC kanallari kisa bir ornekleme ile merkezlenir ve guvenli araliklar otomatik ayarlanir.
\begin{itemize}
  \item Boot sirasinda \term{RC_CALIBRATED}, \term{RC_CALIBRATED_RELAXED} veya gerekirse \term{RC_CALIBRATED_FALLBACK} satiri gorulebilir.
  \item Ardindan \term{RC_CHANNEL_MAP STEERING=CH1 THROTTLE=CH3 ENABLE=CH2} yazisi gelir.
  \item RC merkezleri boot sirasinda olculur.
  \item RC hazir durumu icin once kumandanin notr oldugu dogrulanir.
  \item Ardindan \term{RC_SAFE_READY} satiri gorunur.
  \item CH2 zaten aktifse ve kumanda notrde kalirsa yaklasik \term{800 ms} sonra \term{RC_ENABLE_HELD_AUTO_ARM} ve ardindan \term{RC_ARMED} gorulebilir.
  \item Bu asamada \term{RC_NEUTRAL_WINDOWS} satiri, notr pencerelerin o anki degerini bildirir.
\end{itemize}

\section{Komutlar}
\subsection{Gecerli Komut Formatlari}
Asagidaki komutlar kullanilabilir:
\begin{verbatim}
PWM 100 100
100 100
100,100
100
FL 100
FR 100
RL 100
RR 100
STOP
\end{verbatim}

\subsection{Komutlarin Anlami}
\begin{tabularx}{\linewidth}{@{}p{4.0cm}Y@{}}
\toprule
\textbf{Komut} & \textbf{Anlam} \\
\midrule
\cmd{PWM 100 100} & \cmd{100 100} ile ayni; ileri \\
\cmd{100 100} & Ileri \\
\cmd{100,100} & \cmd{100 100} ile ayni; ileri \\
\cmd{100} & Sol ve sag tarafi ayni degerle ileri sur \\
\cmd{-100 -100} & Geri \\
\cmd{100 -100} & Saga donus / sag pivot \\
\cmd{-100 100} & Sola donus / sol pivot \\
\cmd{FL 100} & Sadece FL teker icin hedef RPM testi \\
\cmd{FR 100} & Sadece FR teker icin hedef RPM testi \\
\cmd{RL 100} & Sadece RL teker icin hedef RPM testi \\
\cmd{RR 100} & Sadece RR teker icin hedef RPM testi \\
\cmd{STOP} & Dur \\
\bottomrule
\end{tabularx}

\subsection{Onemli Not}
\mainsketch{} icinde serial komutlar sadece CH2 pasifken etkili olur.

Ayrica serial pair PWM komutlari tek seferlik degildir; komutun taze kalmasi gerekir. Yaklasik \term{400 ms} icinde yeni komut gelmezse sistem beklemeye gecer. Bu nedenle PC tarafindan \cmd{100 100} gibi pair komutlari ile surus verilecekse komutlar periyodik olarak tekrar gonderilmelidir.

Tek teker RPM test komutlari (\cmd{FL 100}, \cmd{FR 100} gibi) ise \term{STOP} gelene kadar aktif kalir.

\section{Surekli Gelen Seri Satirlar}
Bu bolumde sadece seri hatta surekli akan satirlar anlatilir. Baslangic sirasinda bir kez gorulen kalibrasyon veya acilis mesajlari burada bilerek verilmemistir.

Bu dosyada surekli gorulebilen satirlar:
\begin{itemize}
  \item \term{JS}
  \item \term{RC}
  \item \term{RCDBG}
  \item \term{WTEST}
  \item \term{MODE}
  \item \term{CMD_DIR}
\end{itemize}

\subsection{Telemetri Basliklari}
\subsubsection{JS}
Format:
\begin{verbatim}
JS seq t_us FLrpm10 FRrpm10 RLrpm10 RRrpm10 pwmL pwmR
   FLcmps10 FRcmps10 RLcmps10 RRcmps10 RobotCmps10
\end{verbatim}

Anlam:
\begin{itemize}
  \item \term{seq}: paket sirasi
  \item \term{t_us}: mikro saniye zaman damgasi
  \item \term{FLrpm10}, \term{FRrpm10}, \term{RLrpm10}, \term{RRrpm10}: encoder RPM x10
  \item \term{pwmL}, \term{pwmR}: uygulanan sol ve sag cikis PWM'i
  \item \term{FLcmps10}, \term{FRcmps10}, \term{RLcmps10}, \term{RRcmps10}: teker lineer hizi, \term{cm/s x10}
  \item \term{RobotCmps10}: dort tekerden uretilen robot lineer hizi, \term{cm/s x10}
\end{itemize}

Ornek:
\begin{verbatim}
JS 152 12345678 520 515 501 509 -180 -180 1035 1025 996 1013 1017
\end{verbatim}

Bu ornekte:
\begin{itemize}
  \item FL = 52.0 RPM
  \item FR = 51.5 RPM
  \item RL = 50.1 RPM
  \item RR = 50.9 RPM
  \item sol ve sag cikis PWM = -180
  \item FL hizi = 103.5 cm/s
  \item FR hizi = 102.5 cm/s
  \item RL hizi = 99.6 cm/s
  \item RR hizi = 101.3 cm/s
  \item robot hizi = 101.7 cm/s
\end{itemize}

\subsubsection{RC}
Format:
\begin{verbatim}
RC ch1_us ch2_us ch3_us ch1_ok ch2_ok ch3_ok
\end{verbatim}

Anlam:
\begin{itemize}
  \item \term{ch1_us}: direksiyon kanali
  \item \term{ch2_us}: enable / mode secim kanali
  \item \term{ch3_us}: gaz kanali
  \item \term{ch1_ok}, \term{ch2_ok}, \term{ch3_ok}: sinyal gecerli mi bilgisi
\end{itemize}

\subsubsection{MODE}
Format:
\begin{verbatim}
MODE RC_KUMANDA_AKTIF
MODE SERI_PWM_AKTIF
MODE SERI_WHEEL_AKTIF
MODE BEKLEME
\end{verbatim}

Anlam:
\begin{itemize}
  \item \cmd{MODE RC_KUMANDA_AKTIF}: CH2 aktif, arm tamam ve kontrol RC kumandada
  \item \cmd{MODE SERI_PWM_AKTIF}: CH2 pasif, kontrol serial pair PWM komutlarinda
  \item \cmd{MODE SERI_WHEEL_AKTIF}: CH2 pasif, tek teker RPM test modu aktif
  \item \cmd{MODE BEKLEME}: aktif komut yok, robot beklemede
\end{itemize}

\subsubsection{RCDBG}
Format:
\begin{verbatim}
RCDBG throttlePwm steeringPwm pwmLCmd pwmRCmd pwmLout pwmRout
      neutralReady controlArmed enableReleased enableLatched
\end{verbatim}

Anlam:
\begin{itemize}
  \item Ilk 6 alan RC ic isleminde olusan PWM degerleridir.
  \item Son 4 alan sirasiyla \term{notr hazir}, \term{arm tamam}, \term{enable once birakildi mi}, \term{enable aktif mi} bilgisini verir.
\end{itemize}

\subsubsection{WTEST}
Format:
\begin{verbatim}
WTEST wheel target_rpm measured_rpm10 applied_pwm
      FLrpm10 FRrpm10 RLrpm10 RRrpm10
\end{verbatim}

Anlam:
\begin{itemize}
  \item \term{wheel}: test edilen teker (\term{FL}, \term{FR}, \term{RL}, \term{RR})
  \item \term{target_rpm}: istenen RPM
  \item \term{measured_rpm10}: secili tekerin olculen RPM x10 degeri
  \item \term{applied_pwm}: secili tekere uygulanan PWM
  \item Son 4 alan: tum tekerlerin encoder RPM x10 degerleri
\end{itemize}

\subsubsection{CMD\_DIR}
Format:
\begin{verbatim}
CMD_DIR FWD
CMD_DIR REV
CMD_DIR TURN_LEFT
CMD_DIR TURN_RIGHT
CMD_DIR STOP
CMD_DIR IDLE
CMD_DIR RC_MODE
CMD_DIR FL
CMD_DIR FR
CMD_DIR RL
CMD_DIR RR
\end{verbatim}

Anlam:
\begin{itemize}
  \item \term{FWD}: serial komutta ileri
  \item \term{REV}: serial komutta geri
  \item \term{TURN_LEFT}: serial komutta sola donus
  \item \term{TURN_RIGHT}: serial komutta saga donus
  \item \term{STOP}: serial komutta durus
  \item \term{IDLE}: aktif serial komut yok
  \item \term{RC_MODE}: su anda kontrol RC kumandada
  \item \term{FL}, \term{FR}, \term{RL}, \term{RR}: tek teker RPM test modunda aktif teker
\end{itemize}

\subsection{Ornek Seri Monitor Ciktilari}

\subsubsection{Bosta}
CH2 pasif ve aktif serial komut yokken tipik gorunum:
\begin{verbatim}
JS 31 4070556 0 0 0 0 0 0 0 0 0 0 0
RC 1476 980 1496 1 1 1
MODE BEKLEME
CMD_DIR IDLE
JS 32 4150476 0 0 0 0 0 0 0 0 0 0 0
\end{verbatim}

\begin{tabularx}{\linewidth}{@{}p{4.2cm}Y@{}}
\toprule
\textbf{Deger} & \textbf{Anlam} \\
\midrule
\cmd{31} & JS paket sira numarasi \\
\cmd{4070556} & Mikro saniye zaman damgasi \\
\cmd{0 0 0 0} & FL, FR, RL, RR teker RPM degerleri sifir \\
\cmd{0 0} & Sol ve sag cikis PWM sifir \\
\cmd{0 0 0 0} & FL, FR, RL, RR teker lineer hiz degerleri sifir \\
\cmd{0} & Robot lineer hizi sifir \\
\cmd{1476} & CH1 direksiyon degeri, merkeze yakin \\
\cmd{980} & CH2 pasif, yani RC modu kapali \\
\cmd{1496} & CH3 gaz degeri, merkeze yakin \\
\cmd{1 1 1} & CH1, CH2, CH3 sinyalleri gecerli \\
\cmd{MODE BEKLEME} & Robot beklemede, aktif surus komutu yok \\
\cmd{CMD_DIR IDLE} & Aktif serial yon komutu yok \\
\bottomrule
\end{tabularx}

\subsubsection{RC ile}
CH2 aktif ve kumandadan surus verilirken tipik gorunum:
\begin{verbatim}
JS 3470 280451816 187 200 200 211 -83 -83 371 398 398 419 397
RC 1492 1984 1668 1 1 1
MODE RC_KUMANDA_AKTIF
CMD_DIR RC_MODE
JS 3471 280531984 195 200 205 197 -83 -83 387 398 409 393 397
\end{verbatim}

\begin{tabularx}{\linewidth}{@{}p{4.2cm}Y@{}}
\toprule
\textbf{Deger} & \textbf{Anlam} \\
\midrule
\cmd{3470} ve \cmd{3471} & JS paket sira numaralari \\
\cmd{280451816} ve \cmd{280531984} & Mikro saniye zaman damgalari \\
\cmd{187 200 200 211} & FL, FR, RL, RR icin sirasiyla 18.7, 20.0, 20.0, 21.1 RPM \\
\cmd{-83 -83} & Sol ve sag tarafa uygulanan cikis PWM \\
\cmd{371 398 398 419} & FL, FR, RL, RR icin sirasiyla 37.1, 39.8, 39.8, 41.9 cm/s \\
\cmd{397} & Robot lineer hizi 39.7 cm/s \\
\cmd{1492} & CH1 direksiyon degeri, merkeze yakin \\
\cmd{1984} & CH2 aktif, yani RC modu acik \\
\cmd{1668} & CH3 gaz aktif, ileri/geri surus komutu var \\
\cmd{1 1 1} & Tum RC kanallari gecerli \\
\cmd{MODE RC_KUMANDA_AKTIF} & Kontrol su anda RC kumandada \\
\cmd{CMD_DIR RC_MODE} & Yon bilgisi serialden degil, RC modundan geliyor \\
\bottomrule
\end{tabularx}

\subsubsection{Serial ile}
CH2 pasif ve ornek komut \cmd{120 120} iken tipik gorunum:
\begin{verbatim}
JS 15 2788368 276 284 290 285 -120 -120 549 565 577 567 565
RC 1484 996 1504 1 1 1
MODE SERI_PWM_AKTIF
CMD_DIR FWD
JS 16 2868836 296 301 306 304 -120 -120 589 599 609 605 601
\end{verbatim}

\begin{tabularx}{\linewidth}{@{}p{4.2cm}Y@{}}
\toprule
\textbf{Deger} & \textbf{Anlam} \\
\midrule
\cmd{15} ve \cmd{16} & JS paket sira numaralari \\
\cmd{2788368} ve \cmd{2868836} & Mikro saniye zaman damgalari \\
\cmd{276 284 290 285} & FL, FR, RL, RR icin sirasiyla 27.6, 28.4, 29.0, 28.5 RPM \\
\cmd{-120 -120} & Sol ve sag tarafa uygulanan cikis PWM \\
\cmd{549 565 577 567} & FL, FR, RL, RR icin sirasiyla 54.9, 56.5, 57.7, 56.7 cm/s \\
\cmd{565} & Robot lineer hizi 56.5 cm/s \\
\cmd{1484} & CH1 direksiyon degeri, merkeze yakin \\
\cmd{996} & CH2 pasif, yani serial mod acik \\
\cmd{1504} & CH3 gaz degeri merkeze yakin, surus serial komuttan geliyor \\
\cmd{1 1 1} & Tum RC kanallari gecerli \\
\cmd{MODE SERI_PWM_AKTIF} & Kontrol su anda serial PWM komutlarinda \\
\cmd{CMD_DIR FWD} & Gonderilen serial komut ileri yonunde \\
\bottomrule
\end{tabularx}

\section{Kullanim}
\subsection{RC ile Kullanmak}
\begin{enumerate}
  \item Karta \mainsketch{} yuklu olsun.
  \item Kumandayi ac.
  \item Gaz ve direksiyonu notrde tut.
  \item CH2 aktif konuma al.
  \item Eger CH2 kart acilisinda zaten aciksa, sistem notrde kalirsa otomatik arm olur; gerekirse logda \term{RC_ENABLE_HELD_AUTO_ARM} ve \term{RC_ARMED} satirlarini bekle.
  \item Robotu CH1 ile yon verip, CH3 ile gaz vererek sur.
  \item Gaz altinda donus daha agresiftir; ic taraf daha cok yavaslar, buyuk steering'de ic teker hafif ters yone bile gidebilir.
\end{enumerate}

\subsection{Serial ile Kullanmak}
\begin{enumerate}
  \item Karta \mainsketch{} yuklu olsun.
  \item Portu \term{115200} baud ile ac.
  \item Line ending olarak \term{Newline} sec.
  \item Pair surus icin komut gonder: \cmd{100 100}, \cmd{100}, \cmd{-100 -100}, \cmd{100 -100}, \cmd{-100 100} veya \term{STOP}
  \item Tek teker RPM testi icin: \cmd{FL 100}, \cmd{FR 100}, \cmd{RL 100}, \cmd{RR 100}
\end{enumerate}

\subsection{Serial Mod Kurali}
\mainsketch{} icin kural:
\begin{itemize}
  \item CH2 aktifse serial komut uygulanmaz.
  \item CH2 pasifse serial komut uygulanir.
\end{itemize}

\section{Son Canli Dogrulama}
Guncel kod ile yapilan son kontrol ve dogrulamalarda su davranislar gozlenmistir:
\begin{itemize}
  \item Pair serial komutlarda \cmd{PWM 100 100} ve \cmd{100} ileri yon mantigi ile calisir.
  \item Tek teker RPM test komutlari \cmd{WTEST} satiri ile birlikte kullanilabilir.
  \item RC tarafinda \term{RC_SAFE_READY} sonrasi gerekirse otomatik arm davranisi vardir.
\end{itemize}

Canli RPM ve yon testleri mekanik durum, besleme gerilimi ve bagli yuk durumuna gore degisebilir; bu nedenle sayisal degerler sahada yeniden olculmelidir.

\section{Derleme Dogrulamasi}
Bu manualde anlatilan guncel RC davranisi asagidaki derlemelerle kontrol edilmistir:
\begin{itemize}
  \item \term{4_4_rc_kontrol} \term{arduino:avr:mega} ile derlendi.
  \item \term{4_4_stabil} \term{arduino:avr:mega} ile derlendi.
\end{itemize}

\section{Kisa Notlar}
\begin{itemize}
  \item Yumusak durus vardir; gaz birakildiginda PWM kademeli olarak sifira iner.
  \item Teker hiz dengeleme encoder geri beslemesi ile calisir.
  \item Gaz altinda donus agresif miks ile yapilir; ic taraf daha sert duser ve buyuk steering'de hafif ters yone gecis olabilir.
  \item Ileri giderken bir anda geri komut gelirse, sistem once kontrollu durur sonra ters yone gecer.
  \item Yon degisimi korumasi ortalama hiz ile degil, her taraftaki en hizli tekerin RPM degeri ile calisir. Yani herhangi bir teker hala hizliysa ters yone hemen gecilmez.
  \item Ters yone gecis korumasinda hiz dusumu encoderdan izlenir; ilgili taraf yeterince yavaslamadan yeni yone gecilmez.
  \item Pair serial komut tazeligi yaklasik \term{400 ms} ile izlenir; yeni komut gelmezse sistem beklemeye gecer.
  \item Tek teker RPM test komutlari \term{STOP} gelene kadar aktif kalir.
  \item Kart acilisinda RC tarafi hemen aktif olmaz; once kumandanin notr oldugu gorulmeli, sonra \term{RC_SAFE_READY} durumuna gecilir.
  \item CH2 zaten aktifse ve sistem notrde ise otomatik arm davranisi ile \term{RC_ENABLE_HELD_AUTO_ARM} sonrasi \term{RC_ARMED} gorulebilir.
  \item Boot sirasinda \term{RC_CHANNEL_MAP STEERING=CH1 THROTTLE=CH3 ENABLE=CH2} ve \term{RC_NEUTRAL_WINDOWS} satirlari gorunur.
\end{itemize}

\end{document}
